% !TEX root = ../arbitrage.tex
\section{Why Cheap Talk Persists}
\label{sec:cheap-talk}

Section~\ref{sec:bayesian} identifies pooling on anthropomorphic marketing as a Perfect Bayesian Equilibrium under cheap-talk cost. This section gives five structural reasons for the empirical persistence of that equilibrium. Each is matched to a parameter of the model so that the proposals in Section~\ref{sec:mechanism} can be read as targeted cost changes rather than as welfare exhortations.

\emph{Opacity.} The system-opacity parameter $\theta_S$ is private to the firm. Capability claims about a deployed model cannot be falsified by reading model outputs alone, because the same outputs are compatible with very different training procedures, evaluation regimes, and post-deployment monitoring. Opacity is the precondition for the firm's premium: where outsiders could costlessly verify $\theta_F$, the spread between anthropomorphic and deflationary positioning would collapse. This is the lemons logic of \citet{akerlof1970}: under one-sided private information, the market clears at a price that reflects the worst credible quality, except that here the relevant ``quality'' is a governance disposition rather than a physical attribute.

\emph{Verification bandwidth.} The third inequality of the pooling region, $K_{\theta_R}>\pi(\mathsf{L})D_R$, encodes the regulator's inspection cost relative to expected gain. A low-bandwidth regulator faces an absolute cost so high that abstention is dominant under the prior; even a high-bandwidth regulator may be priced out across the full product population. The condition that produces this state is not unique to AI. \citet{pasquale2015} documents an analogous structure in algorithmic finance and consumer scoring: the opacity of the regulated object combines with the bandwidth of the regulator to make inspection a scarce resource allocated by triage rather than by uniform application.

\emph{Asynchronous harms.} The user's payoff term $H_{\theta_U}$ enters discounted. The episodes in which a relational reading proves costly---deprecation, persona shifts, crisis-handling failures---are temporally separated from the initial attribution by intervals over which the user has already accrued $B_{\theta_U}$. Under standard hyperbolic or even exponential discounting, the inequality $B_{\theta_U}-\pi(\mathsf{L})H_{\theta_U}>0$ holds for a wider class of users than the time-zero risk calculation would suggest. Asynchrony is not a behavioural bug; it is the timing structure of the harm itself.

\emph{Fragmentation of affected parties.} Publics, journalists, intellectuals, and affected communities are folded into the noisy channel $z$ rather than modelled as autonomous players. This is not a normative claim that they lack standing; it is an institutional observation. Their costs of coordination are high, their access to inspection technology is uneven, and their incentives to invest in any particular case are individually small. Even when $z$ carries substantial information about $\theta_F$, regulators cannot treat it as binding without a procedural translation that the channel itself does not provide. The fragmentation result is closely related to the cheap-talk multiplicity of \citet{crawford_sobel1982}: with many would-be senders and a single noisy aggregator, the informative subset of equilibria is not selected for.

\emph{No penalty for inconsistency.} The cheap-talk cost structure sets $c_{\theta_F}(a)=c_{\theta_F}(p)=0$ for both governance types and---more importantly for this section---imposes no joint constraint across $m_F$. Marketing copy and policy documentation are produced under different audiences, different review processes, and different liability regimes. Inconsistency between them is not detected as a single offence; it is decomposed into two compliant statements at two different points of contact. The firm's premium $\Delta_F$ is rent earned precisely on this decomposition. Until the signal cost makes inconsistency itself sanctionable---the proposal of Section~\ref{subsec:sanctionable-inconsistency}---there is no parameter under the firm's control whose adjustment would close the spread.

\emph{Pooling is the market's stable terminal state.} These five structural conditions jointly explain why ethical appeals fail to dislodge the equilibrium. They do not fail because AI governance discourse lacks moral seriousness; they fail because moral seriousness is itself a zero-cost signal. Under the cheap-talk cost structure, a declaration of ``responsible AI'' is formally indistinguishable from any other message at $c_{\theta_F}(a)=0$: both governance types can issue it at the same cost, so it contains no type information and moves no posterior. The equilibrium is not a moral failure. It is the unique stable outcome of the incentive structure as currently constituted, even if the outcome is welfare-reducing in aggregate. A firm that unilaterally forgoes the ontological premium $\Delta_F$---that chooses deflationary marketing across all channels and forfeits the user investment it would otherwise attract---loses revenue it cannot recover from customers who face no such constraint. Market selection eliminates non-arbitraging firms not because restraint is penalized directly, but because the competitor who retains $\Delta_F$ captures the market segment that the restrained firm vacates. The implication is precise: \emph{the market does not reward non-exploitation; it eliminates firms that do not arbitrage}. The mechanisms in Section~\ref{sec:mechanism} respond to this implication by changing the cost parameters rather than appealing to the conscience of players who are already acting rationally.

\begin{table}[t]
\centering
\scriptsize
\setlength{\tabcolsep}{3pt}
\renewcommand{\arraystretch}{1.2}
\begin{tabular}{@{}p{2.7cm}p{2.2cm}p{2.7cm}@{}}
\toprule
\textbf{Structural Condition} & \textbf{Model Parameter} & \textbf{Intervention (\S\ref{sec:mechanism})} \\
\midrule
No penalty for inconsistency & $c_{\theta_F}(a)=c_{\theta_F}(p)=0$ & Sanctionable inconsistency \\
Opacity & $\theta_S$ private to firm & Costly third-party audits \\
Verification bandwidth & $K_{\theta_R}>\pi(\mathsf{L})D_R$ & Auditable records \\
Asynchronous harms & $H_{\theta_U}$ enters discounted & Evidentiary lag priced \\
Fragmentation & $z$ noisy, high coord.\ cost & Aggregation cost lowered \\
\bottomrule
\end{tabular}
\caption{Mapping of structural persistence conditions to model parameters and cost-imposition mechanisms.}
\label{tab:condition-map}
\end{table}
